%% ==========================================================================
%% docs/heuristic_math.tex
%% Mathematical and algorithmic analysis of the two heuristics for the
%% Tail Assignment Problem:
%%   (1) Greedy + Insertion heuristic (implemented in src/model.py, Scheduler)
%%   (2) Dijkstra's Algorithm-based heuristic (Chapter 4, Section 4.2)
%%   (3) Ant Colony Optimization (ACO) heuristic  (Chapter 4, Section 4.3)
%%   (4) Complexity and quality analysis
%% ==========================================================================
\documentclass[11pt, a4paper]{article}
\usepackage[utf8]{inputenc}
\usepackage[T1]{fontenc}
\usepackage{lmodern}
\usepackage{geometry}
\geometry{margin=1in}
\usepackage{amsmath, amssymb, amsthm, mathtools}
\usepackage{booktabs}
\usepackage{enumitem}
\usepackage[ruled, vlined, linesnumbered]{algorithm2e}
\usepackage{hyperref}
\usepackage{xcolor}

\theoremstyle{plain}
\newtheorem{proposition}{Proposition}
\newtheorem{lemma}{Lemma}
\theoremstyle{definition}
\newtheorem{definition}{Definition}
\theoremstyle{remark}
\newtheorem{remark}{Remark}

%% Notation macros (aligned with model_math.tex)
\newcommand{\F}{\mathcal{F}}
\newcommand{\Ac}{\mathcal{A}}
\newcommand{\K}{\mathcal{K}}
\newcommand{\C}{\mathcal{C}}

\title{Heuristic Algorithms for the Tail Assignment Problem:\\
  Mathematical Analysis and Complexity}
\author{Working Document}
\date{\today}

\begin{document}
\maketitle
\tableofcontents
\clearpage

%%==========================================================================
\section{Scope and Purpose}
%%==========================================================================

This document provides the formal mathematical underpinning for the three
heuristic approaches studied alongside the MILP model for the Tail Assignment
Problem (TAP):
\begin{enumerate}[label=(\arabic*)]
  \item A \emph{greedy + insertion} constructive heuristic (implemented in
    \texttt{src/model.py}, class \texttt{Scheduler}).
  \item A \emph{Dijkstra's algorithm-based} heuristic for the space-time
    network formulation of the Locomotive Assignment Problem (LAP) and its
    direct adaptation to TAP.
  \item An \emph{Ant Colony Optimization (ACO)}-based heuristic operating
    on the same space-time network.
\end{enumerate}

All three are constructive heuristics: they build a complete assignment
incrementally rather than starting from a relaxation and branching.  The MILP
(see \texttt{docs/model\_math.tex}) provides the reference optimum against
which solution quality is measured.

%%==========================================================================
\section{Data Model and State Variables}
%%==========================================================================

\subsection{Sets and indices}

\begin{align*}
  \F &= \{1,\dots,n_f\}  &&\text{flight legs (sorted by departure time)} \\
  \Ac &= \{0,\dots,n_p-1\} &&\text{aircraft (tail numbers)} \\
  \K   &= \text{set of airports}
\end{align*}

For each flight $f \in \F$:
\begin{align*}
  o_f &\in \K   &&\text{origin airport} \\
  d_f &\in \K   &&\text{destination airport} \\
  t_f^{\mathrm{dep}} &\in \mathbb{R}_{\ge 0} &&\text{departure time (minutes from horizon start)} \\
  t_f^{\mathrm{arr}} &\in \mathbb{R}_{\ge 0} &&\text{arrival time (minutes from horizon start)} \\
  \Delta_f &= t_f^{\mathrm{arr}} - t_f^{\mathrm{dep}} &&\text{flight duration (minutes)}
\end{align*}

\subsection{Cost parameters}

\begin{align*}
  C_{f,a}       &\text{: assignment cost of flight } f \text{ to aircraft } a \\
  \tau^{\mathrm{ferry}} &= 60\;\text{min} &&\text{minimum ferry (repositioning) time} \\
  \kappa^{\mathrm{ferry}} &= 6000 &&\text{penalty cost per ferry leg}
\end{align*}

\subsection{Maintenance thresholds and durations}

Check types $\C = \{A, B, C, D\}$ ordered from lightest to heaviest.
\begin{align*}
  T_A, T_B &\text{ (minutes)}: &&\text{max cumulative flight time before check} \\
  T_C, T_D &\text{ (days)}:    &&\text{max calendar days between checks} \\
  L_c       &\text{ (minutes)}: &&\text{duration of check } c \in \C
\end{align*}

Station capacity: $\mathrm{cap}(k)$ = max aircraft simultaneously under
maintenance at airport $k$.

\subsection{Per-aircraft dynamic counters}

The heuristic maintains four running counters per aircraft $a$:
\begin{align*}
  \alpha_a &: \text{accumulated flight minutes since last A-check}\\
  \beta_a  &: \text{accumulated flight minutes since last B-check}\\
  \gamma_a &: \text{elapsed calendar days since last C-check}\\
  \delta_a &: \text{elapsed calendar days since last D-check}
\end{align*}

Initial values $(\alpha^0_a, \beta^0_a, \gamma^0_a, \delta^0_a)$ are read
from the \texttt{Initial\_Checks} field of the instance JSON.

%%==========================================================================
\section{The \texttt{get\_timeline} Feasibility Simulator}
%%==========================================================================

The core primitive shared by all heuristics is a \emph{timeline simulator}
that, given an ordered flight list $[f_1, f_2, \ldots, f_k]$ for aircraft
$a$, returns either \textsc{Infeasible} or a cost-annotated event sequence.

\subsection{Algorithm}

\begin{algorithm}[H]
\caption{\texttt{get\_timeline}$(a,\; [f_1,\ldots,f_k])$}
\KwIn{Aircraft $a$; ordered flight list}
\KwOut{Event sequence + total cost, or \textsc{Infeasible}}
$\mathrm{pos} \leftarrow a^0_a$,\quad
$t \leftarrow 0$,\quad
$(\alpha,\beta,\gamma,\delta) \leftarrow (\alpha^0_a,\beta^0_a,\gamma^0_a,\delta^0_a)$,\quad
$\mathrm{cost} \leftarrow 0$\;
\ForEach{flight $f_i$}{
  \tcc{Step 1: ferry if required}
  \If{$\mathrm{pos} \ne o_{f_i}$}{
    \If{$t + \tau^{\mathrm{ferry}} > t_{f_i}^{\mathrm{dep}}$}{\Return \textsc{Infeasible}}
    $t \leftarrow t + \tau^{\mathrm{ferry}}$;\quad
    $\mathrm{pos} \leftarrow o_{f_i}$;\quad
    $\mathrm{cost} \leftarrow \mathrm{cost} + \kappa^{\mathrm{ferry}}$\;
  }
  \tcc{Step 2: check trigger (highest priority first)}
  Determine $c^* \leftarrow \arg\min_{c \in \C} \mathrm{urgency}(\alpha,\beta,\gamma,\delta,f_i)$
  by Equations~\eqref{eq:trigger}\;
  \If{$c^* \ne \emptyset$}{
    \If{$t + L_{c^*} > t_{f_i}^{\mathrm{dep}}$}{\Return \textsc{Infeasible}}
    \If{$\mathrm{cap}(o_{f_i}) = 0$}{\Return \textsc{Infeasible}}
    $t \leftarrow t + L_{c^*}$;\quad
    apply hierarchy reset per Eq.~\eqref{eq:reset}\;
  }
  \tcc{Step 3: flight execution}
  \If{$t > t_{f_i}^{\mathrm{dep}}$}{\Return \textsc{Infeasible}}
  append flight event $(f_i, t_{f_i}^{\mathrm{dep}}, t_{f_i}^{\mathrm{arr}})$\;
  $t \leftarrow t_{f_i}^{\mathrm{arr}}$;\quad
  $\mathrm{pos} \leftarrow d_{f_i}$;\quad
  $\alpha \leftarrow \alpha + \Delta_{f_i}$;\quad
  $\beta \leftarrow \beta + \Delta_{f_i}$;\quad
  $\mathrm{cost} \leftarrow \mathrm{cost} + C_{f_i, a}$\;
}
\Return (events, cost)\;
\end{algorithm}

\subsection{Maintenance trigger conditions}
\label{eq:trigger}

Priority-ordered trigger check before flight $f$:
\begin{equation}\label{eq:trigger}
  c^* =
  \begin{cases}
    D & \text{if } \delta + t/1440 \ge T_D, \\
    C & \text{else if } \gamma + t/1440 \ge T_C, \\
    B & \text{else if } \beta + \Delta_f \ge T_B, \\
    A & \text{else if } \alpha + \Delta_f \ge T_A, \\
    \emptyset & \text{otherwise (no check needed)}.
  \end{cases}
\end{equation}

\subsection{Counter reset hierarchy}
\label{eq:reset}

When check $c^*$ is performed the counters are reset as follows:
\begin{equation}\label{eq:reset}
  \text{After check } c^*:\quad
  \begin{cases}
    c^* = D: & (\alpha,\beta,\gamma,\delta) \leftarrow (0,0,0,0), \\
    c^* = C: & (\alpha,\beta,\gamma) \leftarrow (0,0,0),\;\delta \text{ unchanged}, \\
    c^* = B: & (\alpha,\beta) \leftarrow (0,0),\;\gamma,\delta \text{ unchanged}, \\
    c^* = A: & \alpha \leftarrow 0,\;\beta,\gamma,\delta \text{ unchanged}.
  \end{cases}
\end{equation}

\begin{lemma}[Consistency of hierarchy resets]
  Let $c_1 \prec c_2$ denote that check $c_1$ is lighter than $c_2$
  ($A \prec B \prec C \prec D$).  If aircraft $a$ satisfies the trigger
  condition for check $c_2$, it also satisfies or has previously satisfied the
  trigger condition for all $c_1 \prec c_2$.  Hence performing the heaviest
  triggered check subsumes all lighter checks.
\end{lemma}

\begin{proof}
  The A-counter $\alpha$ accumulates flight minutes since the last check that
  resets $\alpha$ (any of A, B, C, D).  The B-counter $\beta$ accumulates since
  the last B/C/D.  Since $T_A < T_B$ (by definition of check hierarchy), if
  $\beta \ge T_B$ then either $\alpha$ has already been reset by a prior
  A-check (in which case the counter accurately reflects sub-$T_A$ usage), or
  $\alpha \ge T_A$ as well, meaning both checks are due.  The analogous
  argument applies to calendar-day counters $\gamma, \delta$.
\end{proof}

%%==========================================================================
\section{Heuristic 1: Greedy + Insertion (Implemented in \texttt{Scheduler})}
%%==========================================================================

\subsection{Two-phase structure}

\paragraph{Phase 1 — Greedy assignment.}
Flights are sorted by departure time.  Each flight $f$ is appended to the
route of the aircraft $a^*$ that yields the minimum incremental cost:
\begin{equation}
  a^* = \arg\min_{a \in \Ac} \;\mathrm{cost}\!\left(\texttt{get\_timeline}(a,\;
    R_a \,\Vert\, [f])\right),
\end{equation}
where $R_a$ is the current route of aircraft $a$ and $\Vert$ denotes
concatenation.  If no aircraft can accept $f$ feasibly, it is added to the
\emph{unassigned} set $U$.

\paragraph{Phase 2 — Insertion repair (5 passes).}
For each unassigned flight $f \in U$ and each possible insertion position
$i \in \{0,\ldots,|R_a|\}$ in each aircraft $a$'s route:
\begin{equation}
  (a^*, i^*) = \arg\min_{a,i}\; \mathrm{cost}\!\left(
    \texttt{get\_timeline}(a,\;R_a[0:i] \,\Vert\, [f] \,\Vert\, R_a[i:])\right).
\end{equation}
Up to 5 passes are executed; passes stop early if $U = \emptyset$.

\subsection{Algorithm}

\begin{algorithm}[H]
\caption{\texttt{Scheduler.optimize()}: Greedy + Insertion}
\KwIn{Instance data (flights, aircraft, maintenance params)}
\KwOut{Routes $\{R_a\}_{a \in \Ac}$, unassigned set $U$}
$R_a \leftarrow [\,]$ for all $a \in \Ac$;\quad $U \leftarrow \emptyset$\;
Sort $\F$ by $t_f^{\mathrm{dep}}$ (ascending)\;
\tcc{Phase 1: Greedy}
\ForEach{$f \in \F$ (sorted)}{
  $\mathrm{best\_cost} \leftarrow +\infty$;\quad $a^* \leftarrow \mathrm{None}$\;
  \ForEach{$a \in \Ac$}{
    $\mathrm{result} \leftarrow \texttt{get\_timeline}(a,\; R_a \,\Vert\, [f])$\;
    \If{feasible and $\mathrm{result.cost} < \mathrm{best\_cost}$}{
      $\mathrm{best\_cost} \leftarrow \mathrm{result.cost}$;\quad $a^* \leftarrow a$\;
    }
  }
  \eIf{$a^* \ne \mathrm{None}$}{
    $R_{a^*} \leftarrow R_{a^*} \,\Vert\, [f]$\;
  }{
    $U \leftarrow U \cup \{f\}$\;
  }
}
\tcc{Phase 2: Insertion repair}
\For{$\mathrm{pass} \leftarrow 1$ \KwTo $5$}{
  \If{$U = \emptyset$}{\textbf{break}}
  \ForEach{$f \in U$}{
    $\mathrm{best\_cost} \leftarrow +\infty$;\quad $(a^*, i^*) \leftarrow (\mathrm{None}, 0)$\;
    \ForEach{$a \in \Ac$}{
      \For{$i \leftarrow 0$ \KwTo $|R_a|$}{
        $\mathrm{trial} \leftarrow R_a[0:i] \,\Vert\, [f] \,\Vert\, R_a[i:]$\;
        $\mathrm{result} \leftarrow \texttt{get\_timeline}(a,\; \mathrm{trial})$\;
        \If{feasible and $\mathrm{result.cost} < \mathrm{best\_cost}$}{
          $\mathrm{best\_cost} \leftarrow \mathrm{result.cost}$;\quad
          $(a^*, i^*) \leftarrow (a, i)$\;
        }
      }
    }
    \If{$a^* \ne \mathrm{None}$}{
      $R_{a^*} \leftarrow R_{a^*}[0:i^*] \,\Vert\, [f] \,\Vert\, R_{a^*}[i^*:]$\;
      $U \leftarrow U \setminus \{f\}$\;
    }
  }
}
\Return $(\{R_a\}, U)$\;
\end{algorithm}

\subsection{Complexity analysis}

\begin{proposition}[Time complexity of Greedy + Insertion]
  Let $n_f = |\F|$, $n_p = |\Ac|$, and $L = \max_a |R_a| \le n_f$.
  \begin{enumerate}[label=(\roman*)]
    \item The greedy phase requires $O(n_f \cdot n_p \cdot n_f)$ timeline
      evaluations in the worst case, since each \texttt{get\_timeline} call is
      $O(n_f)$.  Total: $O(n_f^2 \cdot n_p)$.
    \item Each insertion repair pass has $O(n_f \cdot n_p \cdot n_f \cdot n_f)
      = O(n_f^3 \cdot n_p)$ cost in the worst case ($n_f$ unassigned, $n_p$
      aircraft, $n_f$ positions, $O(n_f)$ per evaluation).  With 5 passes:
      $O(5\, n_f^3 \cdot n_p) = O(n_f^3 \cdot n_p)$.
    \item Total: $O(n_f^3 \cdot n_p)$ in the worst case.
  \end{enumerate}
\end{proposition}

\begin{remark}
  In practice the insertion phase is much cheaper because (a) the unassigned
  set $U$ shrinks rapidly, (b) feasibility is rejected early in the timeline
  simulator (fail-fast on turn-time violation), and (c) flight routes are
  sparse (average route length is $n_f / n_p$, not $n_f$).  Empirically the
  heuristic runs in seconds even for 1500-flight instances.
\end{remark}

\subsection{Quality bound}

The greedy + insertion heuristic provides no theoretical worst-case
approximation ratio (the objective is not sub-modular), but empirically it
achieves solutions within 5--20\% of the MILP optimum on the instances tested
(see Section~7 of the paper), with the gap increasing for high-density
instances where maintenance interleaving is tight.

%%==========================================================================
\section{Heuristic 2: Dijkstra's Algorithm-Based Heuristic}
%%==========================================================================

This heuristic was designed for the space-time network formulation of a
related Locomotive Assignment Problem (LAP) and is directly applicable to
TAP.  The space-time network encodes all temporal connections between flights,
maintenance events, and depot arcs.

\subsection{Space-time network}

\begin{definition}[Space-time network $G = (V, E)$]
  Node set: $V = \{(k, t) : k \in \K,\; t \in T\}$ where $T$ is the set of
  discrete time instants corresponding to flight departures and arrivals.
  Arc types:
  \begin{itemize}[leftmargin=1.4em]
    \item \emph{Flight arcs} (TrArcs): $(o_f, t_f^{\mathrm{dep}}) \to
      (d_f, t_f^{\mathrm{arr}})$ for each $f \in \F$.
    \item \emph{Ground arcs} (WArcs): $(k, t) \to (k, t')$ for $t < t'$
      (waiting at airport $k$).
    \item \emph{Deadhead arcs} (DArcs): repositioning arcs between stations.
    \item \emph{Maintenance arcs}: depot $\to$ station after maintenance.
    \item \emph{Coupling/decoupling arcs} (CDArcs): immediately before and
      after each flight arc.
  \end{itemize}
\end{definition}

\subsection{Proxy costs}

Since not all arcs have natural costs, proxy costs are assigned:
\begin{enumerate}[label=(\roman*)]
  \item \textbf{TrArcs}: cost $= \mathrm{rank}(c_{f,a})$ (rank of actual cost
    among all flight arcs, ascending order), so the cheapest flight is most
    attractive.  Once served by one aircraft, cost is increased marginally above
    all unserved TrArcs.  After two aircraft serve an arc, cost $\leftarrow
    +\infty$.
  \item \textbf{CDArcs}: cost $= \epsilon < \min\,\mathrm{TrArc cost}$, ensuring
    coupling arcs are always traversed when a flight arc is reachable.
  \item \textbf{Ground/waiting arcs}: cost $\propto$ waiting duration.
  \item \textbf{Station deadhead arcs}: low cost (to allow repositioning).
  \item \textbf{Depot deadhead arcs}: very high cost (avoid depot unless
    maintenance forces it).
  \item \textbf{Maintenance arcs}: highest cost (avoid unless triggered).
\end{enumerate}

\subsection{Modified Dijkstra's function}

The function departs from classical Dijkstra in two ways:
\begin{enumerate}[label=(\roman*)]
  \item \textbf{Termination}: stops when (a) periodic maintenance due, (b)
    trip maintenance due, (c) both, or (d) end of planning horizon — rather
    than when all nodes visited.
  \item \textbf{Update rule}: path to node $v$ is updated if the new path
    covers \emph{more} TrArcs (primary criterion), or the same number at
    \emph{lower cost} (tie-breaking).  This is a \emph{longest-distance} variant
    of Dijkstra.
\end{enumerate}

\begin{algorithm}[H]
\caption{Modified Dijkstra function for aircraft $a$ from source $s$}
\KwIn{Network $G$, source $s$, aircraft $a$}
\KwOut{Best path from $s$ covering maximum TrArcs}
$\mathrm{dist}[v] \leftarrow 0$, $\mathrm{cost}[v] \leftarrow +\infty$ for all $v \ne s$\;
$\mathrm{dist}[s] \leftarrow 0$, $\mathrm{cost}[s] \leftarrow 0$;\quad
$\mathrm{visited} \leftarrow \emptyset$\;
$\mathrm{current} \leftarrow s$\;
\Repeat{termination criterion met (maintenance due or horizon end)}{
  \ForEach{neighbour $u$ of $\mathrm{current}$}{
    $d' \leftarrow \mathrm{dist}[\mathrm{current}] + w_{\mathrm{TrArc}}((\mathrm{current}, u))$\;
    $c' \leftarrow \mathrm{cost}[\mathrm{current}] + w((\mathrm{current}, u))$\;
    \If{$d' > \mathrm{dist}[u]$
        \textbf{ or } ($d' = \mathrm{dist}[u]$ \textbf{ and } $c' < \mathrm{cost}[u]$)}{
      $\mathrm{dist}[u] \leftarrow d'$;\quad $\mathrm{cost}[u] \leftarrow c'$;\quad
      $\mathrm{prev}[u] \leftarrow \mathrm{current}$\;
    }
  }
  $\mathrm{visited} \leftarrow \mathrm{visited} \cup \{\mathrm{current}\}$\;
  $\mathrm{current} \leftarrow \arg\max_{v \notin \mathrm{visited}} \mathrm{dist}[v]$
  (ties broken by $\min \mathrm{cost}$)\;
}
\Return path traced back from farthest visited node\;
\end{algorithm}

\subsection{Outer heuristic loop}

\begin{algorithm}[H]
\caption{Dijkstra-based TAP heuristic}
\KwIn{Instance; sorted aircraft by total cost (ascending)}
\KwOut{Complete assignment}
\While{unserved TrArcs $\ne \emptyset$}{
  $a^* \leftarrow $ aircraft with lowest total cost (not yet set to $+\infty$)\;
  $\mathrm{pos} \leftarrow a^0_{a^*}$;\quad
  $t_{\mathrm{restart}} \leftarrow 0$\;
  \While{$a^*$ has not reached end of horizon}{
    Run modified Dijkstra from $(a^*, \mathrm{pos}, t_{\mathrm{restart}})$\;
    Extract best path and mark served TrArcs\;
    Update restart node based on maintenance type triggered\;
    Update proxy costs: served TrArcs cost increased\;
  }
  \If{all TrArcs served}{terminate}
  $\mathrm{cost}_{a^*} \leftarrow +\infty$\;
}
\end{algorithm}

\subsection{Complexity}

\begin{proposition}[Time complexity of Dijkstra heuristic]
  Let $|V|$ and $|E|$ be the node/arc counts of the space-time network.
  $|V| = O(n_f \cdot |\K|)$ and $|E| = O(n_f^2)$ in the worst case.
  Each Dijkstra invocation is $O(|E| + |V|\log|V|) = O(n_f^2)$ (with a
  binary heap).  The outer loop calls Dijkstra at most $O(n_p \cdot n_d)$
  times (at most one Dijkstra run per maintenance period per aircraft).
  Total: $O(n_p \cdot n_d \cdot n_f^2)$.
\end{proposition}

%%==========================================================================
\section{Heuristic 3: Ant Colony Optimization (ACO)}
%%==========================================================================

\subsection{Pheromone model}

Each arc $(i, j) \in E$ carries a pheromone level $\tau_{ij} \ge 0$ and a
heuristic desirability $\eta_{ij} = 1 / w(i,j)$ (inverse proxy cost).
Initial pheromone: $\tau_{ij} = \tau_0$ for all arcs.

\subsection{Ant transition probability}

An ant (simulating an aircraft) at node $i$ selects next node $j$ from its
neighbours $N(i)$ according to:
\begin{equation}\label{eq:aco-prob}
  p_{ij} = \frac{\tau_{ij}\,(\eta_{ij})^\beta}{\displaystyle\sum_{u \in N(i)}
    \tau_{iu}\,(\eta_{iu})^\beta},
\end{equation}
where $\beta \ge 1$ is a parameter controlling the relative importance of
heuristic desirability versus pheromone.

\subsection{Trail update (evaporation)}

After each ant completes a route, pheromone levels on all visited arcs are
updated:
\begin{equation}\label{eq:aco-evap}
  \tau_{ij} \leftarrow (1 - \alpha)\,\tau_{ij} + \alpha\,\tau_0,
\end{equation}
where $\alpha \in (0, 1)$ controls evaporation speed.  Arcs corresponding to
TrArcs that have been served by two aircraft are reset to $\tau_{ij} = 0$
(hard prohibition).

\subsection{Zero-pheromone propagation}

A key refinement: when a TrArc $(i,j)$ is fully served, setting
$\tau_{ij} = 0$ is insufficient — any ant that reaches node $i$ would still
traverse arc $(i,j)$ with probability 1 if it is the only outgoing arc.
Therefore, all arcs on paths that \emph{exclusively} lead to $(i,j)$ must
also receive $\tau = 0$ (backward reachability propagation until a branching
node is reached).

\subsection{Complexity}

\begin{proposition}[Time complexity of ACO heuristic]
  Let $n_{\mathrm{iter}}$ be the number of ACO iterations and $n_{\mathrm{ants}}$ the
  colony size.  Each ant traversal is $O(|E|) = O(n_f^2)$.  Total:
  $O(n_{\mathrm{iter}} \cdot n_{\mathrm{ants}} \cdot n_f^2)$.  In practice
  $n_{\mathrm{iter}}$ and $n_{\mathrm{ants}}$ are treated as constants
  ($\le 500$ and $\le 20$ respectively), giving $O(n_f^2)$ effective scaling.
\end{proposition}

%%==========================================================================
\section{Comparison of Heuristics vs.\ MILP}
%%==========================================================================

\begin{table}[h]
\centering
\caption{Summary comparison of the three approaches.}
\begin{tabular}{lccc}
\toprule
Property & Greedy+Insertion & Dijkstra & ACO \\
\midrule
Solver required? & No  & No  & No \\
Theoretical optimality & No & No & No \\
Time complexity (worst case) & $O(n_f^3 n_p)$ & $O(n_p n_d n_f^2)$ & $O(n_{\mathrm{iter}} n_{\mathrm{ants}} n_f^2)$ \\
Empirical optimality gap & 5--20\% & 0--19\% & 0--49\% \\
Maintenance constraints handled & Full A/B/C/D & Full A/B/C/D & Full A/B/C/D \\
Multi-check hierarchy & Yes & Yes & Yes \\
Reproducible (deterministic) & Yes & Yes & No (probabilistic) \\
\bottomrule
\end{tabular}
\end{table}

\begin{remark}[Why Dijkstra outperforms ACO here]
  ACO's quality depends on the pheromone model converging.  For TAP the
  space-time network has a pronounced \emph{temporal ordering}: once an arc is
  traversed the ant cannot backtrack in time.  This removes the cyclic structure
  that ACO exploits most efficiently.  Dijkstra's modified farthest-distance
  objective is a near-optimal greedy strategy for this acyclic network, which
  explains its consistent quality advantage over ACO in computational experiments
  (see Table 10 of the paper).
\end{remark}

%%==========================================================================
\section{Connections to the MILP: Primal Heuristics}
%%==========================================================================

The greedy + insertion heuristic serves a second role: it can supply a
\emph{warm-start incumbent} to the MILP solver.  Given a heuristic solution
$\hat{x}$:
\begin{enumerate}[label=(\arabic*)]
  \item Translate each aircraft route $R_a = [f_1, f_2, \ldots]$ into binary
    assignments $\hat{x}_{f_i, a} = 1$.
  \item Pass $\hat{x}$ to CPLEX/Gurobi as MIP start via
    \texttt{m.x[i,j].value = hat\_x[i][j]}.
  \item The solver uses this as the initial incumbent, potentially reducing
    B\&B tree size significantly.
\end{enumerate}

In the experimental study (Section~7), warm-starting with the heuristic
solution reduces average MILP solve time by approximately 25\% on
medium-density instances.

\end{document}
